% =====================================================================
% Section: Certification tables  -- COMPLETE DRAFT
%
% PRIMARY SOURCES for every number in this section:
%   docs/CERTIFICATION_TABLES_2026-07-20.md  (authoritative narrative)
%   results/certification_tables.csv         (machine-readable, 3 margins/family)
%   docs/ATLAS_MINING_REGISTRATION_2026-07-15.md §5 (registered metric set)
%
% DISCREPANCY (surfaced to the author, unresolved):
% docs/CERTIFICATION_TABLES_2026-07-20.md line 4 describes the artifact as
% "12 benchmark families x 3 margins". results/certification_tables.csv contains
% ELEVEN named families (arc_challenge, bbh, gpqa, gsm8k, hellaswag, ifeval,
% math, mmlu, mmlu_pro, musr, winogrande) plus an "ALL (pooled)" row, i.e. 12
% table ROWS at each margin. Table~\ref{tab:certification} below reproduces the
% document's table exactly (11 families + pooled) and the prose says "eleven
% families". Reconcile the doc's line 4 wording before submission.
% =====================================================================

\section{Certification tables: how many items an equivalence claim needs}

\revtwoBanner{every \texttt{required\_n} in the 12-row table, every discordance
quartile, the naive independent-binomial column, the paired-advantage ratios
(1.7$\times$--14.7$\times$, 4.4$\times$ pooled), the per-family cell counts, and
the worked MMLU example. The \emph{method} is unaffected: the estimator, the
one-sided $z$, and the 2\,pp margin do not depend on which cells are in the
atlas.}
\label{sec:certification}

\subsection{Detection is not certification}

The statistical question behind ``near-lossless'' is not the one usually
answered. A McNemar test---including the one-sided variant offered by the
closest existing tool for LLM accuracy statistics
\citep{llmaccuracystats2026}---asks whether there is evidence that the
compressed model differs from its baseline. Failing to find such evidence is not
evidence of equivalence; with a small enough evaluation, nothing is detectable.
Our registration commits to this in advance: ``we will not interpret failure to
reject a difference as equivalence''.
% SOURCE: PREREGISTRATION.md §"Outcomes and Analysis".

Certification asks the other question: \emph{is the difference provably smaller
than a margin I declare in advance?} The standard instrument is TOST---two
one-sided tests at level $\alpha$ each, rejecting the composite null
$|\Delta| \geq m$ in favour of equivalence within $\pm m$. TOST converts an
equivalence claim into something with a sample-size requirement, and this
section computes that requirement empirically.

\subsection{Method}
\label{sec:cert:method}

% SOURCE: docs/CERTIFICATION_TABLES_2026-07-20.md §Method.
For a discordance rate $p_d$ and margin $m$, at 95\% confidence and 80\% power,
\begin{align}
\mathrm{sd}_{\mathrm{paired}}      &= \sqrt{p_d}, &
\mathrm{sd}_{\mathrm{independent}} &= \sqrt{2p(1-p)}, \label{eq:sds}\\[2pt]
n_{\mathrm{req}} &= \left\lceil \left(\frac{(z_{1-\alpha} + z_{1-\beta})\,\mathrm{sd}}{m}\right)^{\!2} \right\rceil
 = \left\lceil \left(\frac{(1.6449 + 0.8416)\,\mathrm{sd}}{m}\right)^{\!2} \right\rceil. \label{eq:nreq}
\end{align}
The paired standard deviation follows from the flip model: the per-item accuracy
difference is $d_i \in \{-1,0,+1\}$, and under the null of no true difference
$\mathrm{Var}(d) = p_d$, the rate at which the two models disagree on
correctness. The independent form uses $p$, the family's median baseline
accuracy, because independent-binomial variance depends on accuracy rather than
on churn. $z_{1-\alpha}$ is \textbf{one-sided}, since TOST rejects two one-sided
nulls at level $\alpha$ each (\S\ref{sec:prereg:choices}, interpretive
choice~5).

The discordance rates are not modelled; they are read off the atlas
(\S\ref{sec:atlas}). For each benchmark family we take the 25th percentile,
median, and 75th percentile of the per-cell accuracy-state churn observed across
that family's atlas cells, giving optimistic, typical, and pessimistic
compression behaviour. Quartiles are \texttt{numpy}'s linear-interpolation
\texttt{np.quantile}.

\subsection{The table}
\label{sec:cert:table}

\begin{table}[t]
\centering
\small
\caption{Items required to certify equivalence within $\pm 2$\,pp at 95\%
confidence and 80\% power, by benchmark family, at the 25th/50th/75th
percentiles of the discordance rates the atlas observes for that family. The
naive column is the same requirement computed by treating the two runs as
independent samples. Eleven benchmark families plus the pooled row.}
\label{tab:certification}
% SOURCE: docs/CERTIFICATION_TABLES_2026-07-20.md §"The table at the registered
% 2 pp margin"; every cell cross-checked against results/certification_tables.csv
% rows with margin_pp = 2.0 (columns n_atlas_cells, discordance_p25/median/p75,
% required_n_p25/median/p75, required_n_independent_binomial,
% paired_advantage_at_median).
\begin{tabular}{lrccrr}
\toprule
Benchmark & Atlas cells & Discordance p25/med/p75 & Required $n$ p25/\textbf{med}/p75 & Naive & Advantage \\
\midrule
musr          &  24 & 0.015 / 0.034 / 0.044 & 232 / \textbf{519} / 681      & 7{,}617 & $14.7\times$ \\
gpqa          &  24 & 0.035 / 0.048 / 0.067 & 545 / \textbf{749} / 1{,}035  & 7{,}233 & $9.7\times$ \\
bbh           & 192 & 0.024 / 0.044 / 0.060 & 371 / \textbf{681} / 928      & 6{,}492 & $9.5\times$ \\
mmlu\_pro     &   5 & 0.048 / 0.053 / 0.059 & 739 / \textbf{827} / 913      & 7{,}695 & $9.3\times$ \\
hellaswag     &  14 & 0.031 / 0.044 / 0.079 & 475 / \textbf{688} / 1{,}222  & 4{,}992 & $7.3\times$ \\
arc\_challenge&   8 & 0.076 / 0.078 / 0.087 & 1{,}172 / \textbf{1{,}211} / 1{,}344 & 7{,}524 & $6.2\times$ \\
ifeval        &   8 & 0.048 / 0.052 / 0.072 & 736 / \textbf{800} / 1{,}108  & 4{,}211 & $5.3\times$ \\
mmlu          & 798 & 0.098 / 0.137 / 0.216 & 1{,}510 / \textbf{2{,}123} / 3{,}343 & 7{,}355 & $3.5\times$ \\
winogrande    &  15 & 0.061 / 0.122 / 0.221 & 940 / \textbf{1{,}879} / 3{,}422 & 6{,}267 & $3.3\times$ \\
math          &  56 & 0.107 / 0.141 / 0.169 & 1{,}661 / \textbf{2{,}186} / 2{,}610 & 5{,}222 & $2.4\times$ \\
gsm8k         &  11 & 0.006 / 0.049 / 0.076 & 100 / \textbf{750} / 1{,}172  & 1{,}236 & $1.7\times$ \\
\midrule
\textbf{ALL (pooled)} & \textbf{1{,}155} & 0.056 / 0.113 / 0.180 & 872 / \textbf{1{,}739} / 2{,}783 & 7{,}663 & $\mathbf{4.4\times}$ \\
\bottomrule
\end{tabular}
\end{table}

\paragraph{How to read a row.} Pick the benchmark family and the equivalence
margin you intend to certify. The three required-$n$ columns are the item counts
you need under optimistic, typical, and pessimistic compression behaviour for
that family. Then compare with the naive column, which is the count you would
compute if you ignored pairing.

% SOURCE: docs/CERTIFICATION_TABLES_2026-07-20.md §"The table" worked example
% and results/certification_tables.csv rows mmlu at margin_pp 1.0 / 2.0 / 3.0
% (required_n_median = 8491 / 2123 / 944; required_n_independent_binomial at
% 2 pp = 7355).
% DISCREPANCY (2026-07-21): docs/CERTIFICATION_TABLES_2026-07-20.md line 33
% states 8,492 for the 1 pp figure; results/certification_tables.csv says 8491.
% The CSV is authoritative -- it is what scripts/certification_tables.py emitted.
% The doc's 8,492 appears to be a naive 4x of the 2 pp value (2,123 x 4), but
% required_n scales as 1/m^2 BEFORE the ceiling: the unrounded 2 pp requirement
% is ~2122.75, so ceil gives 2,123 at 2 pp and 8,491 at 1 pp, not 8,492.
% Text uses 8,491. The doc needs a correction; flagged to Amogh.
\paragraph{Worked example.} \emph{MMLU, 2\,pp margin, typical discordance
$\Rightarrow$ evaluate at least 2{,}123 items.} Ignoring pairing would have
demanded 7{,}355---$3.5\times$ the compute for the same conclusion. Tightening
the margin to 1\,pp raises the requirement to 8{,}491 items; relaxing it to
3\,pp lowers it to 944. The margin is the practitioner's declaration, and its
cost is quadratic; a claim of parity within 1\,pp is four times as expensive to
certify as parity within 2\,pp.

\subsection{Why the naive column belongs in the table}

The independent-binomial column is not a straw man. It is what you get by
treating the baseline and compressed evaluations as two unrelated samples and
comparing proportions, which is the default in most reporting. It is wrong in a
specific, quantifiable way: the two runs are \emph{the same items through two
nearly identical models}, so they agree on the large majority of items, and
their difference has far less variance than independence implies.

% SOURCE: docs/CERTIFICATION_TABLES_2026-07-20.md §"Why the naive column is in
% the table"; results/certification_tables.csv column paired_advantage_at_median.
The advantage column is exactly that gap. It ranges from $\mathbf{1.7\times}$
(GSM8K) to $\mathbf{14.7\times}$ (MuSR) and sits at $\mathbf{4.4\times}$ pooled
over 1{,}155 cells: a practitioner using the paired design reaches the same
equivalence conclusion on roughly a quarter of the evaluation budget. The
variation across families is informative in itself---low-churn families (MuSR,
BBH, GPQA) reward pairing most, while high-churn generative families (MATH,
MMLU) both need more items \emph{and} gain less from pairing.

\subsection{The requirement is driven by churn, not by difficulty}
\label{sec:cert:churn-not-difficulty}

The ordering in Table~\ref{tab:certification} is not the intuitive one, and this
is the section's main conceptual point.

% SOURCE: docs/CERTIFICATION_TABLES_2026-07-20.md, closing paragraph of
% §"Why the naive column is in the table"; results/certification_tables.csv,
% mmlu vs gpqa at margin_pp = 2.0. Baseline accuracies in the next paragraph are
% results/certification_tables.csv column median_baseline_accuracy: gpqa 0.3734,
% mmlu 0.39. Discordances are column discordance_median: mmlu 0.13733,
% gpqa 0.048397.
MMLU needs about \textbf{2{,}123} items at a 2\,pp margin. GPQA needs
\textbf{749}. GPQA is by any ordinary account the harder benchmark---its median
baseline accuracy in the atlas is 0.373 against MMLU's 0.390, on a task designed
to resist exactly the models that saturate MMLU. Difficulty is not what sets the
requirement. What sets it is how much the compressed model's per-item
correctness \emph{churns} relative to the baseline: MMLU's typical discordance
is 0.137 against GPQA's 0.048, so MMLU's paired variance is nearly three times
larger and its required sample size scales with it.

The practical consequence is that a practitioner cannot reason about evaluation
size from intuitions about task hardness, headroom, or score level. Two
benchmarks with similar accuracies and similar apparent difficulty can differ by
a factor of three in the evidence they require, and the only way to know which
is which is to measure churn. That is what the atlas is for, and it is why the
certification tables are empirical rather than analytic.

\subsection{Scope and caveats}
\label{sec:cert:caveats}

% SOURCE: docs/CERTIFICATION_TABLES_2026-07-20.md §"Scope and caveats", all five
% bullets, plus results/certification_tables.csv column n_atlas_cells.
These caveats travel with the table, and any reuse of it should carry them.

\begin{itemize}
  \item \textbf{Families with fewer than 4 analysable cells are omitted.} A
  quartile over three points is not an empirical distribution. This drops
  nothing from the registered set: all surviving families appear in
  Table~\ref{tab:certification}.
  \item \textbf{Four rows rest on thin evidence and are indicative only:}
  \texttt{mmlu\_pro} (5 cells), \texttt{arc\_challenge} (8), \texttt{ifeval}
  (8), and \texttt{gsm8k} (11). By contrast \texttt{mmlu} (798), \texttt{bbh}
  (192), and \texttt{math} (56) are well supported. A reader should treat the
  thin rows as a starting hypothesis to be replaced by their own measured churn,
  not as a reference constant.
  \item \textbf{Family aggregation mixes two sources of variation.} The atlas
  collapses MMLU's 57 per-subject cells, and BBH/MATH/MuSR/GPQA's per-subtask
  cells, into families, so a family's spread mixes subject-level variation with
  model-level variation. This widens the p25--p75 band relative to what a single
  practitioner evaluating a single model would see, which makes the quartile
  columns \emph{conservative} rather than optimistic.
  \item \textbf{Both disclosed feasibility-probe pairs (99 cells) are excluded}
  per the atlas registration §6. They are tiny hand-built sanity pairs---$n$ as
  low as 10, discordance up to 0.9---and would distort every quartile.
  \item \textbf{These are certification sample sizes, not detection sample
  sizes.} They certify equivalence within $\pm m$; the $n$ required to
  \emph{detect} a difference at the same margin is larger. Reporting one as the
  other is the error this section exists to prevent.
\end{itemize}

\subsection{What this section does and does not support}

Supports: an empirically grounded, family-specific answer to ``how many items do
I need to certify parity within $\pm m$?''; a quantified $1.7$--$14.7\times$
($4.4\times$ pooled) reduction in required evaluation from using the paired
design correctly; and the demonstration that the requirement tracks churn rather
than difficulty.

Does not support: extrapolation to benchmark families absent from the atlas;
extrapolation to compression methods absent from the atlas population
(\S\ref{sec:atlas:caveats}); any claim that a model meeting these counts is
equivalent---meeting the count makes the test \emph{informative}, and the test
still has to be run and to pass.
